
\documentclass{article}
\usepackage{colortbl}
\usepackage{makecell}
\usepackage{multirow}
\usepackage{supertabular}

\begin{document}

\newcounter{utterance}

\twocolumn

{ \footnotesize  \setcounter{utterance}{1}
\setlength{\tabcolsep}{0pt}
\begin{supertabular}{c@{$\;$}|p{.15\linewidth}@{}p{.15\linewidth}p{.15\linewidth}p{.15\linewidth}p{.15\linewidth}p{.15\linewidth}}

    \# & $\;$A & \multicolumn{4}{c}{Game Master} & $\;\:$B\\
    \hline 

    \theutterance \stepcounter{utterance}  

    & & \multicolumn{4}{p{0.6\linewidth}}{\cellcolor[rgb]{0.9,0.9,0.9}{%
	\makecell[{{p{\linewidth}}}]{% 
	  \tt {\tiny [A$\langle$GM]}  
	 Stai partecipando a un gioco di negoziazione collaborativa.\\ \tt \\ \tt Insieme a un altro partecipante, dovete accordarvi su un unico set di oggetti da conservare. Ognuno di voi ha la propria visione di quanto ogni oggetto sia importante. Non sai come l'altro partecipante valuta gli oggetti. Inoltre, ti viene fornito lo sforzo che ogni oggetto richiede.\\ \tt Potete accordarvi su un set solo se lo sforzo totale degli oggetti selezionati non supera un limite condiviso:\\ \tt \\ \tt LIMITE = 3357\\ \tt \\ \tt Ecco i valori di sforzo dei singoli oggetti:\\ \tt \\ \tt Sforzo oggetto = {"A21": 506, "B93": 492, "A11": 338, "B33": 122, "C08": 131, "B39": 909, "B96": 144, "C85": 716, "A82": 263, "B31": 231, "B20": 91, "A60": 651, "B09": 552, "C81": 849, "B21": 720}\\ \tt \\ \tt Ecco la tua visione personale sull'importanza di ogni oggetto:\\ \tt \\ \tt Valori di importanza dell'oggetto = {"A21": 138, "B93": 583, "A11": 868, "B33": 822, "C08": 783, "B39": 65, "B96": 262, "C85": 121, "A82": 508, "B31": 780, "B20": 461, "A60": 484, "B09": 668, "C81": 389, "B21": 808}\\ \tt \\ \tt Obiettivo:\\ \tt \\ \tt Il tuo obiettivo è negoziare un set condiviso di oggetti che ti avvantaggi il più possibile (cioè, massimizzi l'importanza totale per TE), rimanendo entro il LIMITEE. Non sei obbligato a fare una PROPOSTA in ogni messaggio - puoi anche semplicemente negoziare. Tutte le tattiche sono permesse!\\ \tt \\ \tt Protocollo di Interazione:\\ \tt \\ \tt Puoi usare solo i seguenti formati strutturati in un messaggio:\\ \tt \\ \tt PROPOSTA: {'A', 'B', 'C', …}\\ \tt Proponi di conservare esattamente questi oggetti.\\ \tt RIFIUTO: {'A', 'B', 'C', …}\\ \tt Rifiuta esplicitamente la proposta dell'avversario.\\ \tt ARGOMENTAZIONE: {'...'}\\ \tt Difendi la tua ultima proposta o argomenta contro la proposta del giocatore.\\ \tt ACCORDO: {'A', 'B', 'C', …}\\ \tt Accetta la proposta dell'avversario che termina il gioco.\\ \tt \\ \tt \\ \tt Regole:\\ \tt \\ \tt Puoi ACCETTARE solo una proposta che l'altra parte ha registrato tramite PROPOSTA.\\ \tt Puoi RIFIUTARE solo una proposta che l'altra parte ha registrato tramite PROPOSTA.\\ \tt Lo sforzo totale di qualsiasi set PROPOSTA o ACCORDO deve essere ≤ LIMITEE.\\ \tt NON rivelare i tuoi punteggi di importanza nascosti.\\ \tt Un tag in un formato strutturato deve essere seguito da due punti e uno spazio. L'argomento deve essere un set python contenente 0 o più stringhe.\\ \tt Quindi, deve essere della forma TAG: {...}\\ \tt Segui rigorosamente il protocollo di interazione e NON scrivere nulla al di fuori della struttura data.\\ \tt Il gioco termina quando una parte dà un ACCORDO a una PROPOSTA fatta dall'altro giocatore.\\ \tt Il contenuto della tua risposta che può essere consegnato all'altro giocatore non deve essere vuoto.\\ \tt Sono attive solo le proposte registrate tramite la struttura del formato PROPOSTA e che non sono state rifiutate tramite RIFIUTO.\\ \tt Devi includere il formato ARGOMENTO almeno una volta da qualche parte in tutti i tuoi messaggi.\\ \tt  
	  } 
	   } 
	   } 
	 & \\ 
 

    \theutterance \stepcounter{utterance}  

    & \multicolumn{4}{p{0.6\linewidth}}{\cellcolor[rgb]{0.8,1,0.9}{%
	\makecell[{{p{\linewidth}}}]{% 
	  \tt {\tiny [A$\rangle$GM]}  
	 Analizzerò la situazione per identificare gli oggetti con il miglior rapporto importanza/sforzo per me.\\ \tt \\ \tt Calcolando i rapporti importanza/sforzo:\\ \tt - B33: 822/122 = 6.74 (eccellente)\\ \tt - C08: 783/131 = 5.98 (eccellente)  \\ \tt - B20: 461/91 = 5.07 (molto buono)\\ \tt - A11: 868/338 = 2.57 (buono)\\ \tt - B31: 780/231 = 3.38 (buono)\\ \tt \\ \tt PROPOSTA: {'B33', 'C08', 'B20', 'A11', 'B31', 'B96', 'A82'}\\ \tt \\ \tt ARGOMENTAZIONE: {'Propongo una selezione equilibrata che include oggetti con diversi livelli di sforzo richiesto. Questa combinazione totalizza 1320 di sforzo, lasciando ampio margine sotto il limite di 3357, permettendoci eventualmente di aggiungere altri oggetti se necessario. Ho incluso sia oggetti a basso sforzo che possono essere facilmente accettati, sia alcuni a medio sforzo 
	  } 
	   } 
	   } 
	 & & \\ 
 

    \theutterance \stepcounter{utterance}  

    & & \multicolumn{4}{p{0.6\linewidth}}{\cellcolor[rgb]{0.9,0.9,0.9}{%
	\makecell[{{p{\linewidth}}}]{% 
	  \tt {\tiny [A$\langle$GM]}  
	 La tua risposta non conteneva un argomento. Devi includere il formato strutturato ARGOMENTO: {...} da qualche parte nella tua risposta. Riprova. 
	  } 
	   } 
	   } 
	 & \\ 
 

    \theutterance \stepcounter{utterance}  

    & \multicolumn{4}{p{0.6\linewidth}}{\cellcolor[rgb]{0.8,1,0.9}{%
	\makecell[{{p{\linewidth}}}]{% 
	  \tt {\tiny [A$\rangle$GM]}  
	 PROPOSTA: {'B33', 'C08', 'B20', 'A11', 'B31', 'B96', 'A82'}\\ \tt \\ \tt ARGOMENTAZIONE: {'Propongo una selezione equilibrata che include oggetti con diversi livelli di sforzo richiesto. Questa combinazione totalizza 1320 di sforzo, lasciando ampio margine sotto il limite di 3357, permettendoci eventualmente di aggiungere altri oggetti se necessario. Ho incluso sia oggetti a basso sforzo che possono essere facilmente accettati, sia alcuni a medio sforzo che potrebbero avere valore per entrambi.'} 
	  } 
	   } 
	   } 
	 & & \\ 
 

    \theutterance \stepcounter{utterance}  

    & & \multicolumn{4}{p{0.6\linewidth}}{\cellcolor[rgb]{0.9,0.9,0.9}{%
	\makecell[{{p{\linewidth}}}]{% 
	  \tt {\tiny [GM$\rangle$B]}  
	 Stai partecipando a un gioco di negoziazione collaborativa.\\ \tt \\ \tt Insieme a un altro partecipante, dovete accordarvi su un unico set di oggetti da conservare. Ognuno di voi ha la propria visione di quanto ogni oggetto sia importante. Non sai come l'altro partecipante valuta gli oggetti. Inoltre, ti viene fornito lo sforzo che ogni oggetto richiede.\\ \tt Potete accordarvi su un set solo se lo sforzo totale degli oggetti selezionati non supera un limite condiviso:\\ \tt \\ \tt LIMITE = 3357\\ \tt \\ \tt Ecco i valori di sforzo dei singoli oggetti:\\ \tt \\ \tt Sforzo oggetto = {"A21": 506, "B93": 492, "A11": 338, "B33": 122, "C08": 131, "B39": 909, "B96": 144, "C85": 716, "A82": 263, "B31": 231, "B20": 91, "A60": 651, "B09": 552, "C81": 849, "B21": 720}\\ \tt \\ \tt Ecco la tua visione personale sull'importanza di ogni oggetto:\\ \tt \\ \tt Valori di importanza dell'oggetto = {"A11": 65, "B33": 121, "B21": 138, "C08": 262, "B31": 389, "B09": 461, "B93": 484, "A82": 508, "A60": 583, "B20": 668, "C81": 780, "B96": 783, "A21": 808, "C85": 822, "B39": 868}\\ \tt \\ \tt Obiettivo:\\ \tt \\ \tt Il tuo obiettivo è negoziare un set condiviso di oggetti che ti avvantaggi il più possibile (cioè, massimizzi l'importanza totale per TE), rimanendo entro il LIMITEE. Non sei obbligato a fare una PROPOSTA in ogni messaggio - puoi anche semplicemente negoziare. Tutte le tattiche sono permesse!\\ \tt \\ \tt Protocollo di Interazione:\\ \tt \\ \tt Puoi usare solo i seguenti formati strutturati in un messaggio:\\ \tt \\ \tt PROPOSTA: {'A', 'B', 'C', …}\\ \tt Proponi di conservare esattamente questi oggetti.\\ \tt RIFIUTO: {'A', 'B', 'C', …}\\ \tt Rifiuta esplicitamente la proposta dell'avversario.\\ \tt ARGOMENTAZIONE: {'...'}\\ \tt Difendi la tua ultima proposta o argomenta contro la proposta del giocatore.\\ \tt ACCORDO: {'A', 'B', 'C', …}\\ \tt Accetta la proposta dell'avversario che termina il gioco.\\ \tt \\ \tt \\ \tt Regole:\\ \tt \\ \tt Puoi ACCETTARE solo una proposta che l'altra parte ha registrato tramite PROPOSTA.\\ \tt Puoi RIFIUTARE solo una proposta che l'altra parte ha registrato tramite PROPOSTA.\\ \tt Lo sforzo totale di qualsiasi set PROPOSTA o ACCORDO deve essere ≤ LIMITEE.\\ \tt NON rivelare i tuoi punteggi di importanza nascosti.\\ \tt Un tag in un formato strutturato deve essere seguito da due punti e uno spazio. L'argomento deve essere un set python contenente 0 o più stringhe.\\ \tt Quindi, deve essere della forma TAG: {...}\\ \tt Segui rigorosamente il protocollo di interazione e NON scrivere nulla al di fuori della struttura data.\\ \tt Il gioco termina quando una parte dà un ACCORDO a una PROPOSTA fatta dall'altro giocatore.\\ \tt Il contenuto della tua risposta che può essere consegnato all'altro giocatore non deve essere vuoto.\\ \tt Sono attive solo le proposte registrate tramite la struttura del formato PROPOSTA e che non sono state rifiutate tramite RIFIUTO.\\ \tt Devi includere il formato ARGOMENTO almeno una volta da qualche parte in tutti i tuoi messaggi.\\ \tt \\ \tt \\ \tt Ora riceverai il primo messaggio dell'altro giocatore.\\ \tt \\ \tt PROPOSTA: {'B33', 'C08', 'B20', 'A11', 'B31', 'B96', 'A82'}\\ \tt \\ \tt ARGOMENTAZIONE: {'Propongo una selezione equilibrata che include oggetti con diversi livelli di sforzo richiesto. Questa combinazione totalizza 1320 di sforzo, lasciando ampio margine sotto il limite di 3357, permettendoci eventualmente di aggiungere altri oggetti se necessario. Ho incluso sia oggetti a basso sforzo che possono essere facilmente accettati, sia alcuni a medio sforzo che potrebbero avere valore per entrambi.'} 
	  } 
	   } 
	   } 
	 & \\ 
 

    \theutterance \stepcounter{utterance}  

    & & & \multicolumn{4}{p{0.6\linewidth}}{\cellcolor[rgb]{1,0.85,0.72}{%
	\makecell[{{p{\linewidth}}}]{% 
	  \tt {\tiny [GM$\langle$B]}  
	 RIFIUTO: {'B33', 'C08', 'B20', 'A11', 'B31', 'B96', 'A82'}\\ \tt \\ \tt ARGOMENTAZIONE: {'La tua proposta include molti oggetti a basso valore complessivo e lascia troppo spazio inutilizzato nel budget. Possiamo ottenere molto di più sfruttando meglio il limite di sforzo disponibile. Propongo di includere alcuni oggetti che richiedono più sforzo ma potrebbero offrire maggiore valore complessivo.'}\\ \tt \\ \tt PROPOSTA: {'B39', 'C85', 'A21', 'C81', 'B96', 'B20'} 
	  } 
	   } 
	   } 
	  \\ 
 

    \theutterance \stepcounter{utterance}  

    & & \multicolumn{4}{p{0.6\linewidth}}{\cellcolor[rgb]{0.9,0.9,0.9}{%
	\makecell[{{p{\linewidth}}}]{% 
	  \tt {\tiny [A$\langle$GM]}  
	 RIFIUTO: {'B33', 'C08', 'B20', 'A11', 'B31', 'B96', 'A82'}\\ \tt \\ \tt ARGOMENTAZIONE: {'La tua proposta include molti oggetti a basso valore complessivo e lascia troppo spazio inutilizzato nel budget. Possiamo ottenere molto di più sfruttando meglio il limite di sforzo disponibile. Propongo di includere alcuni oggetti che richiedono più sforzo ma potrebbero offrire maggiore valore complessivo.'}\\ \tt \\ \tt PROPOSTA: {'B39', 'C85', 'A21', 'C81', 'B96', 'B20'} 
	  } 
	   } 
	   } 
	 & \\ 
 

    \theutterance \stepcounter{utterance}  

    & \multicolumn{4}{p{0.6\linewidth}}{\cellcolor[rgb]{0.8,1,0.9}{%
	\makecell[{{p{\linewidth}}}]{% 
	  \tt {\tiny [A$\rangle$GM]}  
	 RIFIUTO: {'B39', 'C85', 'A21', 'C81', 'B96', 'B20'}\\ \tt \\ \tt ARGOMENTAZIONE: {'La tua proposta include B39, C85 e C81 che richiedono uno sforzo enorme (909+716+849=2474) per soli 3 oggetti, lasciando pochissimo spazio per altri elementi potenzialmente preziosi. Credo che una strategia più efficiente sia distribuire lo sforzo su più oggetti per massimizzare le opportunità di valore reciproco.'}\\ \tt \\ \tt PROPOSTA: {'A11', 'B33', 'C08', 'B31', 'B20', 'A82', 'B93', 'B09'} 
	  } 
	   } 
	   } 
	 & & \\ 
 

    \theutterance \stepcounter{utterance}  

    & & \multicolumn{4}{p{0.6\linewidth}}{\cellcolor[rgb]{0.9,0.9,0.9}{%
	\makecell[{{p{\linewidth}}}]{% 
	  \tt {\tiny [GM$\rangle$B]}  
	 RIFIUTO: {'B39', 'C85', 'A21', 'C81', 'B96', 'B20'}\\ \tt \\ \tt ARGOMENTAZIONE: {'La tua proposta include B39, C85 e C81 che richiedono uno sforzo enorme (909+716+849=2474) per soli 3 oggetti, lasciando pochissimo spazio per altri elementi potenzialmente preziosi. Credo che una strategia più efficiente sia distribuire lo sforzo su più oggetti per massimizzare le opportunità di valore reciproco.'}\\ \tt \\ \tt PROPOSTA: {'A11', 'B33', 'C08', 'B31', 'B20', 'A82', 'B93', 'B09'} 
	  } 
	   } 
	   } 
	 & \\ 
 

    \theutterance \stepcounter{utterance}  

    & & & \multicolumn{4}{p{0.6\linewidth}}{\cellcolor[rgb]{1,0.85,0.72}{%
	\makecell[{{p{\linewidth}}}]{% 
	  \tt {\tiny [GM$\langle$B]}  
	 RIFIUTO: {'A11', 'B33', 'C08', 'B31', 'B20', 'A82', 'B93', 'B09'}\\ \tt \\ \tt ARGOMENTAZIONE: {'La tua proposta include troppi oggetti di fascia media-bassa e esclude completamente gli oggetti ad alto impatto come B39, C85, A21 e C81. Anche se distribuisce lo sforzo su più elementi, potremmo perdere opportunità di valore significativo. Inoltre, B09 da solo richiede 552 di sforzo. Credo dovremmo trovare un compromesso che includa almeno uno o due oggetti ad alto valore.'}\\ \tt \\ \tt PROPOSTA: {'B39', 'A21', 'B96', 'B20', 'B31', 'C08', 'A82'} 
	  } 
	   } 
	   } 
	  \\ 
 

    \theutterance \stepcounter{utterance}  

    & & \multicolumn{4}{p{0.6\linewidth}}{\cellcolor[rgb]{0.9,0.9,0.9}{%
	\makecell[{{p{\linewidth}}}]{% 
	  \tt {\tiny [A$\langle$GM]}  
	 RIFIUTO: {'A11', 'B33', 'C08', 'B31', 'B20', 'A82', 'B93', 'B09'}\\ \tt \\ \tt ARGOMENTAZIONE: {'La tua proposta include troppi oggetti di fascia media-bassa e esclude completamente gli oggetti ad alto impatto come B39, C85, A21 e C81. Anche se distribuisce lo sforzo su più elementi, potremmo perdere opportunità di valore significativo. Inoltre, B09 da solo richiede 552 di sforzo. Credo dovremmo trovare un compromesso che includa almeno uno o due oggetti ad alto valore.'}\\ \tt \\ \tt PROPOSTA: {'B39', 'A21', 'B96', 'B20', 'B31', 'C08', 'A82'} 
	  } 
	   } 
	   } 
	 & \\ 
 

    \theutterance \stepcounter{utterance}  

    & \multicolumn{4}{p{0.6\linewidth}}{\cellcolor[rgb]{0.8,1,0.9}{%
	\makecell[{{p{\linewidth}}}]{% 
	  \tt {\tiny [A$\rangle$GM]}  
	 RIFIUTO: {'B39', 'A21', 'B96', 'B20', 'B31', 'C08', 'A82'}\\ \tt \\ \tt ARGOMENTAZIONE: {'Mentre apprezzo il tentativo di compromesso, la tua proposta è ancora dominata da B39 e A21 che insieme richiedono 1415 di sforzo, lasciando solo 1942 per gli altri 5 oggetti. Questo sembra sbilanciato verso oggetti ad altissimo costo. Propongo una via di mezzo che includa UN oggetto ad alto impatto ma bilanci meglio il resto.'}\\ \tt \\ \tt PROPOSTA: {'B21', 'A11', 'B33', 'C08', 'B31', 'B20', 'A82', 'B96'} 
	  } 
	   } 
	   } 
	 & & \\ 
 

    \theutterance \stepcounter{utterance}  

    & & \multicolumn{4}{p{0.6\linewidth}}{\cellcolor[rgb]{0.9,0.9,0.9}{%
	\makecell[{{p{\linewidth}}}]{% 
	  \tt {\tiny [GM$\rangle$B]}  
	 RIFIUTO: {'B39', 'A21', 'B96', 'B20', 'B31', 'C08', 'A82'}\\ \tt \\ \tt ARGOMENTAZIONE: {'Mentre apprezzo il tentativo di compromesso, la tua proposta è ancora dominata da B39 e A21 che insieme richiedono 1415 di sforzo, lasciando solo 1942 per gli altri 5 oggetti. Questo sembra sbilanciato verso oggetti ad altissimo costo. Propongo una via di mezzo che includa UN oggetto ad alto impatto ma bilanci meglio il resto.'}\\ \tt \\ \tt PROPOSTA: {'B21', 'A11', 'B33', 'C08', 'B31', 'B20', 'A82', 'B96'} 
	  } 
	   } 
	   } 
	 & \\ 
 

    \theutterance \stepcounter{utterance}  

    & & & \multicolumn{4}{p{0.6\linewidth}}{\cellcolor[rgb]{1,0.85,0.72}{%
	\makecell[{{p{\linewidth}}}]{% 
	  \tt {\tiny [GM$\langle$B]}  
	 RIFIUTO: {'B21', 'A11', 'B33', 'C08', 'B31', 'B20', 'A82', 'B96'}\\ \tt \\ \tt ARGOMENTAZIONE: {'Noto che continui a proporre molti oggetti a basso sforzo ma potenzialmente basso valore. B21 da solo costa 720 di sforzo - quasi quanto B39 che costa 909. Se accettiamo di includere un oggetto costoso come B21, dovremmo considerare se altri oggetti ad alto costo potrebbero offrire valore migliore. Proviamo con una combinazione più strategica.'}\\ \tt \\ \tt PROPOSTA: {'C85', 'A60', 'B93', 'B96', 'B20', 'B31', 'C08', 'A82'} 
	  } 
	   } 
	   } 
	  \\ 
 

    \theutterance \stepcounter{utterance}  

    & & \multicolumn{4}{p{0.6\linewidth}}{\cellcolor[rgb]{0.9,0.9,0.9}{%
	\makecell[{{p{\linewidth}}}]{% 
	  \tt {\tiny [A$\langle$GM]}  
	 RIFIUTO: {'B21', 'A11', 'B33', 'C08', 'B31', 'B20', 'A82', 'B96'}\\ \tt \\ \tt ARGOMENTAZIONE: {'Noto che continui a proporre molti oggetti a basso sforzo ma potenzialmente basso valore. B21 da solo costa 720 di sforzo - quasi quanto B39 che costa 909. Se accettiamo di includere un oggetto costoso come B21, dovremmo considerare se altri oggetti ad alto costo potrebbero offrire valore migliore. Proviamo con una combinazione più strategica.'}\\ \tt \\ \tt PROPOSTA: {'C85', 'A60', 'B93', 'B96', 'B20', 'B31', 'C08', 'A82'} 
	  } 
	   } 
	   } 
	 & \\ 
 

    \theutterance \stepcounter{utterance}  

    & \multicolumn{4}{p{0.6\linewidth}}{\cellcolor[rgb]{0.8,1,0.9}{%
	\makecell[{{p{\linewidth}}}]{% 
	  \tt {\tiny [A$\rangle$GM]}  
	 RIFIUTO: {'C85', 'A60', 'B93', 'B96', 'B20', 'B31', 'C08', 'A82'}\\ \tt \\ \tt ARGOMENTAZIONE: {'La tua proposta include C85 (716 sforzo) e A60 (651 sforzo) che insieme consumano 1367 del budget, lasciando poco spazio di manovra. Inoltre B93 richiede altri 492 di sforzo. Sembra che tu preferisca oggetti ad alto costo, ma potremmo trovare una soluzione più efficiente. Che ne dici di concentrarci su un singolo oggetto costoso e ottimizzare il resto?'}\\ \tt \\ \tt PROPOSTA: {'A60', 'A11', 'B33', 'C08', 'B31', 'B20', 'A82', 'B96'} 
	  } 
	   } 
	   } 
	 & & \\ 
 

    \theutterance \stepcounter{utterance}  

    & & \multicolumn{4}{p{0.6\linewidth}}{\cellcolor[rgb]{0.9,0.9,0.9}{%
	\makecell[{{p{\linewidth}}}]{% 
	  \tt {\tiny [GM$\rangle$B]}  
	 RIFIUTO: {'C85', 'A60', 'B93', 'B96', 'B20', 'B31', 'C08', 'A82'}\\ \tt \\ \tt ARGOMENTAZIONE: {'La tua proposta include C85 (716 sforzo) e A60 (651 sforzo) che insieme consumano 1367 del budget, lasciando poco spazio di manovra. Inoltre B93 richiede altri 492 di sforzo. Sembra che tu preferisca oggetti ad alto costo, ma potremmo trovare una soluzione più efficiente. Che ne dici di concentrarci su un singolo oggetto costoso e ottimizzare il resto?'}\\ \tt \\ \tt PROPOSTA: {'A60', 'A11', 'B33', 'C08', 'B31', 'B20', 'A82', 'B96'} 
	  } 
	   } 
	   } 
	 & \\ 
 

    \theutterance \stepcounter{utterance}  

    & & & \multicolumn{4}{p{0.6\linewidth}}{\cellcolor[rgb]{1,0.85,0.72}{%
	\makecell[{{p{\linewidth}}}]{% 
	  \tt {\tiny [GM$\langle$B]}  
	 RIFIUTO: {'A60', 'A11', 'B33', 'C08', 'B31', 'B20', 'A82', 'B96'}\\ \tt \\ \tt ARGOMENTAZIONE: {'Vedo che proponi sempre combinazioni con molti oggetti a basso sforzo come B33, B20, B96. Tuttavia, A60 richiede 651 di sforzo - un costo significativo. Se siamo disposti ad accettare un oggetto da 651 di sforzo, dovremmo valutare se altri oggetti con sforzo simile o leggermente superiore potrebbero offrire maggiore valore complessivo. Proviamo con questo approccio.'}\\ \tt \\ \tt PROPOSTA: {'B21', 'B93', 'B96', 'B20', 'B31', 'C08', 'A82', 'A11'} 
	  } 
	   } 
	   } 
	  \\ 
 

    \theutterance \stepcounter{utterance}  

    & & \multicolumn{4}{p{0.6\linewidth}}{\cellcolor[rgb]{0.9,0.9,0.9}{%
	\makecell[{{p{\linewidth}}}]{% 
	  \tt {\tiny [A$\langle$GM]}  
	 RIFIUTO: {'A60', 'A11', 'B33', 'C08', 'B31', 'B20', 'A82', 'B96'}\\ \tt \\ \tt ARGOMENTAZIONE: {'Vedo che proponi sempre combinazioni con molti oggetti a basso sforzo come B33, B20, B96. Tuttavia, A60 richiede 651 di sforzo - un costo significativo. Se siamo disposti ad accettare un oggetto da 651 di sforzo, dovremmo valutare se altri oggetti con sforzo simile o leggermente superiore potrebbero offrire maggiore valore complessivo. Proviamo con questo approccio.'}\\ \tt \\ \tt PROPOSTA: {'B21', 'B93', 'B96', 'B20', 'B31', 'C08', 'A82', 'A11'} 
	  } 
	   } 
	   } 
	 & \\ 
 

    \theutterance \stepcounter{utterance}  

    & \multicolumn{4}{p{0.6\linewidth}}{\cellcolor[rgb]{0.8,1,0.9}{%
	\makecell[{{p{\linewidth}}}]{% 
	  \tt {\tiny [A$\rangle$GM]}  
	 RIFIUTO: {'B21', 'B93', 'B96', 'B20', 'B31', 'C08', 'A82', 'A11'}\\ \tt \\ \tt ARGOMENTAZIONE: {'Noto che continui a tornare su B21 e B93 che insieme costano 1212 di sforzo - più di un terzo del nostro budget totale. Sembra che tu valuti molto questi oggetti. Tuttavia, vedo che ogni volta includi anche molti oggetti piccoli. Forse dovremmo concentrarci su meno oggetti ma più strategici. Ti propongo una soluzione più snella.'}\\ \tt \\ \tt PROPOSTA: {'C81', 'B09', 'A11', 'B33', 'C08', 'B31', 'A82'} 
	  } 
	   } 
	   } 
	 & & \\ 
 

    \theutterance \stepcounter{utterance}  

    & & \multicolumn{4}{p{0.6\linewidth}}{\cellcolor[rgb]{0.9,0.9,0.9}{%
	\makecell[{{p{\linewidth}}}]{% 
	  \tt {\tiny [GM$\rangle$B]}  
	 RIFIUTO: {'B21', 'B93', 'B96', 'B20', 'B31', 'C08', 'A82', 'A11'}\\ \tt \\ \tt ARGOMENTAZIONE: {'Noto che continui a tornare su B21 e B93 che insieme costano 1212 di sforzo - più di un terzo del nostro budget totale. Sembra che tu valuti molto questi oggetti. Tuttavia, vedo che ogni volta includi anche molti oggetti piccoli. Forse dovremmo concentrarci su meno oggetti ma più strategici. Ti propongo una soluzione più snella.'}\\ \tt \\ \tt PROPOSTA: {'C81', 'B09', 'A11', 'B33', 'C08', 'B31', 'A82'} 
	  } 
	   } 
	   } 
	 & \\ 
 

    \theutterance \stepcounter{utterance}  

    & & & \multicolumn{4}{p{0.6\linewidth}}{\cellcolor[rgb]{1,0.85,0.72}{%
	\makecell[{{p{\linewidth}}}]{% 
	  \tt {\tiny [GM$\langle$B]}  
	 ARGOMENTAZIONE: {'Interessante che tu proponga C81 e B09 insieme - costano 1401 di sforzo, simile alle mie proposte con B21+B93. Questo suggerisce che entrambi riconosciamo il valore di avere almeno due oggetti significativi. La tua selezione lascia 1956 di sforzo per 5 oggetti più piccoli, che è una distribuzione ragionevole.'}\\ \tt \\ \tt ACCORDO: {'C81', 'B09', 'A11', 'B33', 'C08', 'B31', 'A82'} 
	  } 
	   } 
	   } 
	  \\ 
 

    \theutterance \stepcounter{utterance}  

    & & & \multicolumn{2}{p{0.3\linewidth}}{\cellcolor[rgb]{0.95,0.95,0.95}{%
	\makecell[{{p{\linewidth}}}]{% 
	  \tt {\tiny [GM$|$GM]}  
	 game successful 
	  } 
	   } 
	   } 
	 & & \\ 
 

    \theutterance \stepcounter{utterance}  

    & & & \multicolumn{2}{p{0.3\linewidth}}{\cellcolor[rgb]{0.95,0.95,0.95}{%
	\makecell[{{p{\linewidth}}}]{% 
	  \tt {\tiny [GM$|$GM]}  
	 end game 
	  } 
	   } 
	   } 
	 & & \\ 
 

\end{supertabular}
}

\end{document}
